\section{Introduction}
\label{sec:intro}

% World models are becoming a trending topic recently. Currently, people are trying to have world model to be grounded by physics laws, as the neural network are more likely to learn the geometrical correlation relationships rather than actual causal relationship between the change of appearance and movement. Thus, it's being an active research topic in the world model community to somehow having physics prior as an intermediate layer of the world model. However, even when the physics is correct, there's still gap between the world model and the reality, which is the behavioral gap. The driving behavior of one typical human driver on a highway covered by snow is inherently different from the one in the sunny day. The world model ought to understand that difference and reflect it accurately. In order to achieve the behavioral realism, we attempt by having the world model video generation based on the behavior of vehicles trained by real motion data. Note that in autonomous driving study, there exists tremendous amount of data, despite the point of view being in the dashcam view, but not in the bird eye view (BEV). Thus, the video generation of those foundation models are naturally biased towards car view. With that as a constraint, it's necessary to introduce more guidance for the current world foundation models to generate accurate videos. 

According to Federal Highway Administration (FHWA) \cite{noauthor_how_nodate}, there occurs 6,000,000 vehicle crashes each year, 12 percent of the crashes are related to weather. Out of all the adverse weather related crashes,  94 percent of the car crashes are related to the change of pavement friction, including 73 percent of the crash happening during rain or mist conditions, 21 percent happening when there's frozen precipitation.

Human driving behavior naturally changes when driving under different weather conditions. Depending on the weather, the visibility of the road, and the dynamic reaction of the car changes, and human drivers tend to capture those changes and react accordingly. For example, human drives more cautiously when the rainfall is high, as the dual effect of occluded visibility and the loss of maneuverability as the effective friction of the road decreases. However, most existing world models and simulation frameworks do not explicitly capture this behavioral adaptation, which may lead to unrealistic driving behavior in simulation. Accurate world responses are important for applications such as transportation planning, traffic safety analysis, and risk assessment of road systems under different weather conditions, etc. 

Therefore, an important question arises: how can we construct a world model that both follows physical laws and reproduces realistic driving behavior? This problem consists of two closely related components that are causally dependent: driving behavior depends on the underlying world physics. Driving behavior has been widely studied, and there exist multiple structured datasets that record expert driving trajectories. However, because people are less likely to travel during extreme weather conditions, the data distribution in these datasets is often biased toward more common weather scenarios. As a result, learning driving behavior purely from data may fail to generalize to rare environmental conditions. In addition, many open driving datasets contain limited or coarse-grained weather annotations, which further complicates learning weather-dependent driving behavior.

This work proposes a pipeline for training weather-aware driving agents, using real human driving data to guide the training process. The trained agents are then used as actors to generate traffic scenes that are both physically realistic and behaviorally plausible. The paper is organized into three parts. The first part introduces the process of converting weather conditions into friction parameters that can be consumed by the simulation environment. The second part describes the training design of the pipeline. The final part discusses the integration of the generated traffic scenes into the overall simulation framework.
[More will be written in the introduction]
% THIS PARAGRAPH CAN DEFINITELY BE BETTER   
% This work features the integration of a weather to road condition module, that takes water film depth and road materials as inputs, convert them to effective static and dynamic friction, that is applied to the ground of a world in Isaac Sim. The world is later used to host agents to learn driving behaviors. After the agnets learned how to drive, they are deployed into real scenes extracted from Waymo Open Motion Dataset. A world that is generated from the waymo road geometry will be simulated, waypoints generated for agents, and the simulation is ran. A camera will be established at a fixed location, and the traffic scene will be rendered from isaacsim. FInally, the captured video will have its edge extracted, and being used as an input of Cosmos transfer, for video generation.

%-------------------------------------------------------------------------
